\section{Experimental Setup}
\label{sec:setup}

\paragraph{Architectures and model scales.} 
We instantiate DeepSeek-V4.1-Flash's asymmetric CED topology across an $L{=}\val{cfg/layers}$-layer stack (Figure~\ref{fig:main}a), split evenly into an encoder (layers 0--5) and a decoder (layers 6--11). Every layer maintains a local SWA cache over a $w{=}\val{cfg/window}$-token window. In the encoder, layers 0--1 use SWA only, while layer 2 serves as the sole encoder global KV source, read by layers 3--5. At the boundary, layer 6 projects the shared decoder global KV from the encoder output, read by layers 7--11. All global-reading layers merge local SWA and global KV in a single softmax (\S\ref{sec:background}). Each token thus stores only two global KV entries—one from the encoder and one from the decoder. 
% To isolate CED's cross-layer KV topology, we use standard dense attention and FFNs without orthogonal compression or MoE modules. 
To isolate CED's cross-layer KV topology, we use standard multi-head attention and dense FFNs without intra-layer KV compression (e.g., GQA or MLA) or MoE modules.
Our dense Transformer baseline uses the exact same dimensions but equips every layer with full-causal KV and no SWA, caching \val{kvratio/x40/1024/dense-over-ced}$\times$ more global KV per token than CED. Fixing this 12-layer layout, we scale hidden dimensions across $d \in \{768, 1024, 1536\}$, yielding approximately 90M, 150M, and 340M non-embedding parameters over a \val{cfg/seq}-token context (adding a loop introduces $\sim$1\% parameters for input injection; full architectural specifications are in Appendix~\ref{app:hparams}).

\paragraph{Compared configurations and ablations.} 
At each scale, we compare an unlooped baseline against a 4-layer recurrent loop ($T{=}\val{cfg/rounds}$ rounds) placed at layers 2--5 for both CED (\emph{encoder loop}) and the dense model, alongside a CED \emph{decoder loop} at layers 7--10. At 90M and 150M, we also train a dense model looped at the middle third (layers 4--7) to match Parcae's default placement~\citep{parcae}. Finally, two 150M CED ablations isolate our core mechanisms: (i) a \emph{fixed-depth} encoder loop trained with a deterministic \val{cfg/rounds}-round unroll to test whether cross-round KV sharing requires stochastic-depth training (\S\ref{sec:res-kv}); and (ii) a \emph{YOCO-U} variant~\citep{yocou} that loops layers 2--5 within an SWA-only encoder (making layer 6 the sole global KV source), testing whether the loop itself must recompute global KV (\S\ref{sec:res-longdoc}).

\paragraph{Training and evaluation.} 
Models are trained from scratch on FineWeb-Edu~\citep{fineweb} using Muon~\citep{moonlight_muon} and AdamW under a unified learning-rate schedule (Appendix~\ref{app:hparams}). Standard runs allocate $\sim$\val{cfg/ratio/x40/1024} tokens per non-embedding parameter (\val{cfg/tokens/x40/768}B, \val{cfg/tokens/x40/1024}B, and \val{cfg/tokens/x40/1536}B tokens for 90M, 150M, and 340M). To evaluate extended scaling, we also train the 150M models for \val{cfg/tokens/long}B tokens, matching the budget of 370M Parcae~\citep{parcae}, our 150M configuration (Parcae counts embeddings). Within each scale and budget, all runs see identical data orders, matching parameter counts and training tokens rather than FLOPs. 

We report upstream bits per byte (BPB) on FineWeb-Edu validation and downstream zero/few-shot accuracy on DCLM CORE and CORE-Extended~\citep{dclm}. Furthermore, because dependencies beyond CED's \val{cfg/window}-token SWA window rely exclusively on global KV—and standard FineWeb-Edu documents are too short to span our \val{cfg/seq}-token context (median \val{data/doc-median} tokens)—we evaluate position-wise BPB across contiguous \val{cfg/seq}-token chunks drawn from three long-document corpora: PG19~\citep{pg19}, FinePDFs ($\ge 4{,}096$ tokens)~\citep{kydlicek2025finepdfs}, and GovReport ($\ge 2{,}048$ tokens)~\citep{govreport} (protocol details in Appendix~\ref{app:longctx}).
